%!TeX root=csp-blueprint.tex
\section{The Lean development}\label{sec:formalization}

\subsection{What is reused}\label{sec:reuse}

The predecessor development \texttt{ZhukLean} is a complete, \texttt{sorry}-free
Lean~4 formalization of Zhuk's centre theorem, built on
\texttt{Mathlib.ModelTheory}. It is consumed here as a \texttt{lake} dependency,
not copied. It supplies:

\begin{itemize}\tightlist
\item products of structures (\texttt{piStructure}, \texttt{prodStructure},
  coordinatewise realization, reindexing and evaluation homomorphisms) --- the
  one part of the universal-algebra vocabulary Mathlib lacks;
\item absorption in the tuple form (\texttt{Witnesses}, \texttt{Absorbs},
  \texttt{BinAbsorbs}), Taylor terms (\texttt{IsTaylorOn}), and central
  absorption (\texttt{CentrallyAbsorbs});
\item the relational description of absorption by essential relations
  (Barto--Kazda), which is what converts absorption at some arity into
  absorption at arity~$3$;
\item Zhuk's centre theorem itself (\texttt{zhuk\_center}) and the ternary
  collapse (\texttt{exists\_ternary\_\brk witnesses}) --- which is precisely
  \inproof{lem:central-ternary}, cited without proof in \cite{ZhukSimplified}.
\end{itemize}

Those also discharge, in advance, every ingredient of \cite[Lem.~6.11]{ZhukStrong},
which the proof paper's stable-intersection theorem needs in its type-$\TC$
case: essentiality, the Barto--Kazda criterion, and ``central implies ternary
absorbing''.

\subsection{What is built here}\label{sec:built}

\begin{center}
\begin{tabular}{@{}llr@{}}
\hline
Module & Contents & Status \\
\hline
\texttt{CSP/Product} & products of structures & complete \\
\texttt{CSP/Congruence} & \texttt{Congruence L M}, order, \texttt{sInf}, quotient & complete \\
\texttt{CSP/Irreducible} & irreducibility, existence and uniqueness of $\cover\sigma$ & complete \\
\texttt{CSP/Bridge} & bridges, collapse, composition, linear/PC & 2 imports open \\
\texttt{CSP/Types} & the six types, multi-types, $\lll$ & complete \\
\texttt{CSP/Instance} & instances, reductions, consistency, cruciality & complete \\
\texttt{CSP/Main} & the main statements & targets \\
\hline
\end{tabular}
\end{center}

Everything below \texttt{CSP/Main} is \texttt{sorry}-free. The two open items in
\texttt{CSP/Bridge} are \inproof{lem:bridge-comp} and its collapse identity,
both imported from \cite{ZhukJACM}. These are claims about a repository, not
about the mathematics; they should be read against a commit hash, and they are
deliberately kept out of the proof paper for that reason.

\subsection{Design decisions}\label{sec:design}

\begin{description}
\item[Build on \texttt{Mathlib.ModelTheory}, keep the language generic.]
  A bespoke ``finite algebra with one WNU operation'' type would bundle the
  standing hypotheses as fields and avoid threading them, which is a real cost
  in the predecessor development. But it would discard
  \texttt{Substructures.lean} --- 985 lines of exactly the $\Sg$ API needed,
  including \texttt{mem\_closure\_iff\_exists\_term} with the variable type
  equal to the generating set, an interface no hand-rolled clone matches --- and
  it would discard the predecessor's 1600 lines wholesale. Congruences and
  quotients are also \emph{cheaper} in \texttt{ModelTheory}, not dearer.
  Instantiating to $\Vn$ is a thin layer on top. Carry the standing hypotheses
  as typeclass instances on the structure rather than as ambient hypotheses, so
  that the audit of which result consumes which hypothesis is mechanical.

\item[Terms and term operations are different types.]
  \cite{ZhukSimplified} writes $t\in\Clo(\mathbf A)$ and treats $t$ as both a
  syntactic term and the operation it induces. Every statement of the proof
  paper that quantifies over $\Clo(\mathbf A)$ is about \emph{operations}, so
  the syntactic layer may be quotiented away --- but it cannot be skipped,
  because $\Sg$ is characterised through terms and because the arity
  bookkeeping in the absorption arguments is syntactic.

\item[$\lll$ is a witnessed chain, not \texttt{Relation.ReflTransGen}.]
  An earlier draft of this section said the opposite, and it was wrong: the
  proof paper's ``the congruences coming from $C\lll^{A}B$'' is a function of
  the chain and not of the pair, and two arguments turn on \emph{which} chain
  is chosen. So the primary notion is an inductive family carrying, at each
  step, the intermediate set, its type, and its witnessing congruence, with
  extension by one step as an operation on that data;
  \texttt{Relation.ReflTransGen} is its propositional truncation and may be
  used only where no congruence is named. Both are needed: the truncation
  supplies the induction principle most consumers want.

\item[Dotted relations are defined as disjunctions.]
  \texttt{DotLll C B := C = }$\varnothing$\texttt{ }$\vee$\texttt{ Lll C B}, so
  the case split is forced at every use site rather than resolved by the
  reader. Conflating $\lll$ with $\dlll$ is the commonest way to misread the
  type theory of the proof paper's \S3.

\item[Absorption is stated over coordinatewise-constrained tuples.]
  That is: for every $\bar a\in A^{k}$ with $a_{j}\in B$ for all $j\neq i$, one
  has $t(\bar a)\in B$. \emph{Not} as ``for all
  $b_{1},\dots,b_{k}\in B$ and $a\in A$, $t(b_{1},\dots,a,\dots,b_{k})\in B$''.
  The two agree, but the second phrasing
  invites the variant in which $b_{i}$ is quantified and then discarded, which
  is vacuously true when $B=\varnothing$ and at $k=1$ where the intended
  content is not vacuous. This was an actual defect in the predecessor
  blueprint.

\item[The bridge criterion is the definition of linearity.]
  Zhuk defines a linear congruence by a per-block isomorphism to
  $\mathbb Z_{p}^{m}$ and derives the bridge criterion. We reverse this: the
  criterion is first-order over finite data and is what every consumer uses;
  the block description becomes a structure theorem. See
  \inproof{haz:linear-def}.

\item[$\cover\sigma$ is data attached to irreducibility.]
  \texttt{Irreducible.exists\_isCover} produces it; there is no standalone
  function, because there is no cover without irreducibility. A pleasant
  by-product: with the definition taken literally, irreducibility \emph{implies}
  $\sigma\neq A^{2}$, since the empty family intersects to the universe.

\item[$\ConOne$ gets a neutral name until it is a congruence.]
  The raw relation is always reflexive and symmetric but not transitive; the
  notation invites the reader to assume congruence-hood. Reserve $\ConOne$ for
  the case where subdirectness and rectangularity make it one. Note also that
  it depends on the \emph{position} in the scope, not on the variable, so it is
  ill-defined if a variable occurs twice.

\item[Instances are lists; scopes are tuples; the variable type is a parameter.]
  Proofs delete a constraint and reason about $\mathcal I\setminus\{C\}$, and
  the same relation may occur twice on the same tuple inside an expanded
  covering, so \texttt{List} with membership by index is the safe choice.
  Repeated variables in a scope are allowed and occur.

\item[Coverings are built by coproduct of variable types, not by renaming.]
  \cite{ZhukSimplified} writes ``we rename the variables so that the only
  common variable of $\Upsilon_{x}$ and $\mathcal J$ is $x$''. With a concrete
  variable type this means a fresh-name supply, a renaming operation, and a
  lemma that renaming preserves each of $1$-consistency, cycle-consistency,
  cruciality, and being a covering --- eight lemmas the paper does not state.
  Making the variable type a parameter and gluing over $V_{1}\oplus V_{2}$
  quotiented by the shared variable turns all eight into transport along an
  equivalence. This is the single most important representation decision below
  the algebra, and it should be made before anything about coverings is
  written.

\item[Nothing is stubbed with \texttt{sorry} at the level of a definition.]
  \texttt{IsConnectedInstance} and \texttt{IsExpandedCovering} are
  \emph{absent} rather than defined as \texttt{sorry}, because a
  \texttt{def~:~Prop~:=~sorry} is a junk constant that downstream proofs can
  silently lean on. Absence forces the design decision; a stub hides it.
\end{description}

\subsection{What Mathlib is missing}\label{sec:mathlib-gaps}

\begin{center}\small
\begin{tabular}{@{}p{0.32\textwidth}p{0.60\textwidth}@{}}
\hline
\textbf{Layer} & \textbf{Status} \\
\hline
Products of first-order structures &
  Absent from Mathlib; supplied by \texttt{ZhukLean} (\S\ref{sec:reuse}). \\
Congruences of a structure, quotients &
  \texttt{ModelTheory} has substructures but not congruences; built here in
  \texttt{CSP/Congruence}. \\
Absorption, central subuniverses &
  Absent; supplied by \texttt{ZhukLean}. \\
Mixed-prime ``dimension'' &
  A subgroup of $\prod_{i}\mathbb Z_{q_{i}}$ with distinct $q_{i}$ is not a
  vector space. \texttt{ZMod} and \texttt{Module (ZMod p)} exist; the primary
  decomposition into $\mathbb F_{q}$-vector spaces and the codimension-one
  characterisation must be built. Needed by the proof paper's codimension-one
  theorem. \\
The Pol--Inv Galois connection &
  Absent. Roughly four printed pages, of independent value, and the single
  largest item in the \textup{(H0)} budget. \\
Complexity classes, $\le_{p}$, \textup{NP} &
  Absent; see \S\ref{sec:complexity-layer}. Built separately in $\approx 850$
  lines. \\
Polynomial completeness, Istinger--Kaiser &
  Absent. Needed only for \inproof{lem:pc-on-top}, which nothing else consumes;
  a staged development should treat ``PC congruence'' as ``irreducible and not
  linear'' throughout and prove that lemma last or not at all. \\
\hline
\end{tabular}
\end{center}

\subsection{Order of attack}\label{sec:order}

\begin{enumerate}\tightlist
\item \inproof{lem:ba-center-implies}. Small, self-contained, used everywhere,
  and adjacent to what is already formalized. Its statement needs both
  hypotheses the source drops --- subdirectness, and a common binary term ---
  and the common term should be carried in the type, as \cite{ZhukStrong}
  writes $\le_{\TBA(t)}$.
\item The mixed-prime linear algebra of \S\ref{sec:mathlib-gaps}. Independent
  of everything else; can be done in parallel.
\item \inproof{lem:bridge-comp}. Twelve source lines of relational algebra;
  closes the two open items in \texttt{CSP/Bridge}.
\item The coproduct design for coverings, then the basic properties of expanded
  coverings.
\item \inproof{thm:stable-intersection}. The one hard statement in the
  interface, and the sole source of bridges.
\item \inproof{lem:bridge-from-instance}. Eight hypotheses, ten pages.
\item \inproof{thm:main-inductive}. Last.
\end{enumerate}

\subsection{Scope, and what governs the schedule}\label{sec:scope}

The calibration point is that the predecessor formalized about $3.5$ printed
pages in about $1670$ lines of Lean --- roughly $480$ lines per page. Excluding
\S\ref{sec:complexity-layer} and \cite{ZhukSimplified}'s \S4, which is an
independent second result and not needed for the dichotomy, the transitive
closure of what must be proved is on the order of $100$--$130$ printed pages,
so $35\,000$--$60\,000$ lines of Lean.

Line count is not the binding constraint, and it should not be converted into
person-time here: any such figure would be a claim about a development process
rather than about this mathematics, and it is not evidence a reader of this
document can check.

What binds is the \emph{critical path}: the depth of the serial chain, and the
number of places where the source is wrong and new mathematics is needed. The
serial chain is
\begin{center}\small
\pkey{thm:stable-intersection}\\[2pt]
$\longrightarrow$\ \pkey{lem:bridge-from-instance}\\[2pt]
$\longrightarrow$\ \pkey{thm:main-inductive}
\end{center}
each a single large proof that cannot be split across workers --- stable
intersection is the sole source of bridges, bridge-from-instance is the only
consumer that turns them into instance-level structure, and the main induction
consumes that.

The new mathematics is a much shorter list than an earlier draft of this
section claimed, and the items on it are not the ones it named. Of the audit
document's register, thirteen defects are repaired with no new mathematics.
What remains open, and gates the formalization of \S5 of
\cite{ZhukSimplified}, is:

\begin{itemize}\tightlist
\item Claim \pkey{cl:typing-anchor} and \inproof{lem:typed-above} --- the
  residue of the typing half of \inproof{lem:parallelogram-crucial}. They
  replaced \inproof{lem:maximal-mult} and its unproved hypothesis on the
  critical path; see the audit document. Nothing downstream of them should
  be attempted first.
\item The rectangularity normalisation ``compose the bridge with itself enough
  times'', on which \inproof{lem:no-cross-bridge} and
  \inproof{lem:bridge-to-pc} rest. The reflexivisation that used to stand
  beside it on this list is \emph{not} open: it is an observation about a
  hypothesis every consumer already carries, and it is now
  \inproof{lem:bridge-reflexive}.
\item \inproof{lem:multiply-all-linear}, which is not yet a proof: its
  Claim \pkey{cl:multiply-step} is open.
  (\inproof{lem:self-intersection-pc} and \inproof{lem:intersection-good},
  which stood here, are now both proved in full.)
\end{itemize}

Those are not throughput-limited. Schedule around them, not around the line
count.
